\cotefile{64}
\scanpage{53}\mspage{5}
\[
 a)\longrightarrow b)\qquad \Gamma=\underline{\mathrm{Norm}}_G\,\underline V
\]
$b)\longrightarrow a)$ \quad $\underline V$ est l'unique sous-schéma en groupes \inserted{distingué} fini étale de $\Gamma$, \deleted{\unsure{g}de} degré $4$, \deleted{distingué}.
\margin{\uline{NB} $\Gamma/V\simeq\underline{\mathrm{Aut}}_{S\mathrm{gr}}(V)\simeq\underline{\mathrm{Aut}}_S(V^*)$}
\[
 c)\longrightarrow b)\qquad \Gamma=\underline{\mathrm{Aut}}_S(X,\Delta)\xleftarrow{\sim}\underline{\mathrm{Aut}}^+_{S\text{-oct}}(\Delta)
\]
$a)\longrightarrow c)$ \quad \deleted{$\underline V^*$ est} Toute section $\sigma$ de $V^*=V\smallsetminus(\text{section unité})$ définit une action \inserted{fidèle} de \deleted{\unsure{ce|st}} $(\mathbb Z/2\mathbb Z)_S$ sur $X$, donc un ``axe'' $\Delta_\sigma$ de $X$, et $\Delta$ est \deleted{la réunion} (loc. ét.) le \deleted{\unsure{sché}} diviseur somme de ses axes \ldots

\uline{Corollaire} 1 \quad Pour un diviseur étale de degré 6 $\Delta$ dans $X$, il existe au plus une structure icosaédrale régulière~; \deleted{\unsure{dont} le \inserted{groupe} \illegible{1} des autom.} \inserted{S'il en existe une, le} \inserted{schéma en} groupes de ses automorphismes \inserted{directs} est isom. à $\underline{\mathrm{Aut}}_S^{\deleted{\unsure{+|4}}}(X,\Delta)$.

\uline{Corollaire 2} \quad La catégorie des couples $(X,\Delta)$ \inserted{\uline{sur $S$}}, (ou $(X,\underline V)$, $(X,\underline\Gamma)$ \ldots) \deleted{\unsure{orientés}} \inserted{est} équivalente à celle des octaèdres \deleted{\unsure{orientés}} combinatoires \inserted{$\Delta$} \inserted{sur} $S$, ou encore à celle des revêtements finis étales \deleted{$\unsure{X|V|\underline V}$} \inserted{$\widetilde F$} de degré 4 de $S$, en associant à \unsure{tt} $(X,\Delta)$ le revêtement des quatre ``faces modulo antipodisme'' de la structure octaédrale $\Delta$, ou ce qui revient au $=$, celui des \inserted{3-}sous-groupes de Sylow de \deleted{$\unsure{G|\underline{\mathfrak S}}_{\Delta}$} $\underline{\mathrm{Aut}}^+_{S\text{-oct}}(\Delta)\simeq\underline{\mathrm{Aut}}_S(X,\Delta)$.
\inserted{munis d'un iso $\underline{\unsure{\mathrm{Or}|\mathrm{or}}}\,\Delta/S\simeq\mu^*_{4S}$ (on dit qu'ils sont $\mu^*_4$-orientés)}
La structure-type $*$ correspond à \deleted{\unsure{inscrire à|brouillon|revêt.}} $X\subset\mathbb P^2_{\mathbb Z[\frac12]}$, $X=\mathbb V(x^2+y^2+z^2=0)$, avec l'op. de \deleted{$\unsure{V}$} $\{\pm\}^3/\mathrm{diag}\simeq\mathbb V_{\{0,1,\infty\}}(\mathbb F_2)$ donnée par les chgts de signe en $x,y,z$ $*$.
% margin note written sideways
\margin{\uline{NB} Ici $X$ est totalement anisotrope sur $\mathbb R$ i.e. n'a pas de pts réels sur $\mathbb R$ -- c'est le revêtement galoisien \unsure{universel} \unsure{absolument} de groupe $\mathbb V_{\{0,1,\infty\}}(\mathbb F_2)$ de $\mathbb P^1_{\deleted{\unsure{\mathbb Z}}}$, ramifié \unsure{seulem\textsuperscript{t}|seulement} en $0,1,\infty$ \unsure{(en} \ldots}
% margin note written sideways
\margin{$*$ La \inserted{bi}face distinguée est formée des deux faces $\{(0,1,i),(i,0,1),(1,i,0)\}$ et $\{\{0,1,-i\},\{-i,0,1\},\{1,-i,0\}\}$}
% archive sheet number: 52
